\section{Conclusions and Limitations}
\label{subsec:q:scope}

\paragraph{Answer to the research question.}
Yes: the institution does harbour latent intellectual communities
recoverable from publication records alone.  The strongest evidence
is the GW/astrophysics group---33 UniPi researchers whose shared
intellectual context is invisible in the citation graph but fully
legible in shared external references.  More broadly, all 200 gap
candidates localise in BC-corrected communities, concentrate in 6
disciplinarily coherent communities (FOS~L4 purity $\geq 45\%$,
five of six $\geq 60\%$), and are confirmed across three algorithms
and a bootstrap null ($p < 0.0005$).

\paragraph{What the evidence does not show.}

\textit{No ground truth.}  The analysis is entirely structural:
thematic proximity and community coherence are necessary but not
sufficient to confirm a missed collaboration.

\textit{Temporal actionability is limited.}  Sixty-nine per cent of
candidates pair two pre-2020 papers, documenting historical proximity
rather than current opportunity.  The 9~Recent--Recent pairs are the
operationally relevant subset.

\textit{Within-discipline scope.}  The pipeline operates within
FOS~L4 sub-disciplines at the current debiasing threshold;
cross-disciplinary gap detection would require a different scoring
function.

\paragraph{Directions for enhancement.}
The pipeline identifies \emph{where to look}; three targeted
enhancements would convert its output into actionable
recommendations.  \textit{Temporal filtering}---restricting
candidates to Recent--Recent pairs---eliminates the 69\% reflecting
historical accumulation as a one-line filter on existing output.
\textit{Author-level lifting} (nodes = researchers, edges =
co-authorship or co-citation via ORCID) would enable recommendations
of the form ``researcher $X$ and researcher $Y$ share 15 reference
clusters but have never co-authored''---contingent on disambiguation
coverage.  \textit{Retrospective validation}---flagging 2020
candidates and checking for co-authorship by 2025---would yield the
first empirical precision and recall estimates for BC-based
collaboration prediction at the institutional level.  The
7~genuine Recent--Recent pairs identified in
Section~\ref{subsec:q:temporal} are the natural starting point: few
enough for manual review, recent enough to be actionable, and robust
across all three algorithms.

\textit{Multi-hop proximity.}  Both components measure affinity at the
first hop only: shared direct references (HE) or shared direct
citation neighbours (BC).  Extending to second-order
metrics---bibliographic coupling on the hypergraph, personalized
PageRank on the citation graph---would capture transitive proximity
currently invisible to the scoring function.